\section{Conclusions}\label{sec:conclusions}

We formulated stochastic firebreak placement on reduced directed
propagation graphs and characterized scenario recourse by reachability.
Its LP has a componentwise minimal integral burn vector for every
binary placement and relatively complete recourse. Downstream,
capped-coverage and path inequalities give valid master bounds. The
complete path family recovers fractional recourse on general directed
graphs, while coverage coincides with it on rooted trees; the
implemented path separator retains only a limited path collection.

Across 810 matched reduced \texttt{new20x20} runs, Combinatorial
B\&B certifies 95.6\% (43/45) of expected-loss cases and has mean final MIP gap
0.18\% and mean observed solver time 212.4\,s. Path-exp is the
strongest recorded configuration under CVaR (46.7\% certified, mean
gap 3.51\%) and mean--CVaR (66.7\%, mean gap 2.10\%). These are
objective-dependent optimization outcomes at a requested relative
gap $10^{-3}$ and a 1,800\,s stopping limit. Unresolved runs and
recorded time overruns enter the gap and time summaries, so their
mean observed runtime is not unrestricted time to optimality.

A complementary 135-placement Path-exp study evaluates five sampled
decisions for each objective, size and budget on a fixed 1,000-scenario
test set. At $\alpha=0.02$, increasing training size from 100 to 400
reduces mean held-out expected loss from 94.95 to 91.73, CVaR from
204.00 to 192.76, and mean--CVaR from 155.42 to 143.44 burned cells.
Placement similarity varies by objective and budget: for example,
mean pairwise Jaccard similarity under CVaR at $\alpha=0.02$ is
0.08, 0.27 and 0.22 across the three sizes. These comparisons describe
the decisions found by Path-exp. At $\alpha=0.03$ and $n=400$ none
of its CVaR runs and only one of its five mean--CVaR runs is certified,
so unresolved optimization error limits causal interpretation of the
risk-objective trends. Evaluating each selected placement on paired
reduced and full Reburn test graphs would establish whether the
observed reduced-graph improvements transfer to the full model.
